\documentclass[11pt]{article}
\usepackage[hmargin=.94in,vmargin=.78in]{geometry}
\usepackage{amsmath,amssymb}
\usepackage[hidelinks]{hyperref}
\hypersetup{pdftitle={The first Mills number on fourteen vertices},
  pdfauthor={Troy Valiyaparampil and GPT-6.2 Astra}}
\newtheorem{theorem}{Theorem}
\newtheorem{lemma}[theorem]{Lemma}
\newenvironment{proof}{\par\noindent\textit{Proof.}\ }{\hfill\(\square\)\par}
\newcommand{\F}{\mathcal F}
\newcommand{\G}{\mathcal G}
\title{The first Mills number on fourteen vertices}
\author{Troy Valiyaparampil and GPT-6.2 Astra}
\date{25 November 2026}
\begin{document}
\maketitle
\begin{abstract}
In a distribution of data across fourteen machines for which each pair of
items shares a machine, the largest machine load is at least \(7/23\) of
the data. Colin Barker's 2008 covering of 69 points by fourteen blocks of
size 21 attains this fraction. This establishes the \(q=3\) case of
a conjecture posed by F\"uredi in 1990. We prove the matching lower bound for
arbitrary real weights. A count of incidences on four-element patterns
reduces a putative counterexample to twelve such patterns and two
five-element patterns. The intersection of the latter two forces a
packing of ten pairs that can accommodate at most eight.
\end{abstract}

\section{The parameter and the construction}
Let an item have as its \emph{pattern} the set of machines holding it.
Pairwise comparison is possible exactly when every two patterns intersect.
Merge items with the same pattern and give that pattern their combined
weight. On the vertex set \(V=\{0,\ldots,13\}\), let \(w_A\ge0\)
be the weight on each member of an intersecting family of nonempty
subsets. The load at \(x\) is \(\ell_x=\sum_{A\ni x}w_A\).
The infimum of \(\max_x\ell_x\) over all such families with total weight
one is the all-to-all comparison data limit \(L(14)\) of Hall, Horsley and
Stinson \cite{hhs}. This is also the first Mills number \(M_1(14)\) in the
notation of Guiduli and Kir\'aly \cite{gk}. With vertex capacities one,
the largest possible total weight is its reciprocal.

\begin{theorem}\label{thm:main}
\(L(14)=M_1(14)=7/23\). The maximum fractional matching number of
an intersecting hypergraph on fourteen vertices is \(23/7\).
\end{theorem}

F\"uredi \cite[Conjecture 4.9]{furedi94} wrote \(\mu(n)\) for the
maximum fractional matching number of an intersecting hypergraph on at
most \(n\) vertices. Isolated vertices allow us to use exactly \(n\).
In 1990 he conjectured \(\mu(q^2+q+2)\le q+2/(2q+1)\), with equality
when a projective plane of order \(q\) exists \cite[Conjecture (6.3)]{furedi90}.
Theorem~\ref{thm:main} establishes this prediction for \(q=3\).

Barker's covering in the La Jolla archive \cite{barker} already gives
\(L(14)\le21/69=7/23\) by Lemma~3.4 of \cite{hhs}. For an independently
reproducible witness, the table lists the same covering up to relabeling,
recovered by local search. In order, the fourteen blocks in Barker's
archive map to coordinates \(5,1,6,10,7,11,0,9,8,2,4,3,12,13\) in the table.
Each two patterns meet, every coordinate has load 21, and the
weights sum to 69. Taking \(w_A\) distinguishable copies of pattern \(A\)
produces fourteen distinct 21-element blocks covering all
\(\binom{69}{2}=2346\) pairs. The same witness establishes strict
improvement on the projective-plane value, since \(7\cdot13<4\cdot23\).

\begin{center}\small
\begin{tabular}{rl@{\qquad}rl}
Weight & Pattern & Weight & Pattern\\ \hline
5 & \(\{9,10,11,12\}\) & 5 & \(\{2,3,7,10\}\)\\
3 & \(\{0,5,7,11\}\) & 6 & \(\{0,4,8,10,13\}\)\\
5 & \(\{3,4,5,12\}\) & 5 & \(\{1,4,7,9\}\)\\
5 & \(\{6,7,8,12\}\) & 5 & \(\{1,3,8,11\}\)\\
5 & \(\{2,5,8,9\}\) & 6 & \(\{0,3,6,9,13\}\)\\
5 & \(\{2,4,6,11\}\) & 6 & \(\{0,1,2,12,13\}\)\\
5 & \(\{1,5,6,10\}\) & 3 & \(\{5,7,11,13\}\)
\end{tabular}
\end{center}

One can read the construction directly from the projective plane of
order three. Split a point into \(0\) and \(13\). Keep the nine lines
avoiding it, each with weight 5. On one of the four lines through it,
retain the other three points \(\{5,7,11\}\) and adjoin \(0\) or \(13\)
separately to obtain two patterns of weight 3. On each remaining line,
replace the split point by \emph{both} new points and give the resulting
five-set weight 6. Every old point receives \(3\cdot5+6=21\), and each
new point receives \(3+3\cdot6=21\).

\section{A deficit on four-element patterns}
To prove the other direction, suppose all loads are at most 21 and the
total weight \(W\) exceeds 69. Any normalized distribution with maximum
load below \(7/23\) can be scaled to these bounds. On its support maximize
total mass subject to the fourteen capacities and nonnegativity. Each
variable is bounded by 21, so a maximizing extreme point exists. Its
positive incidence columns are linearly independent, since any dependence
allows small perturbations in both directions preserving all loads and
nonnegativity. We may therefore assume at most fourteen positive
patterns. This is the support reduction of \cite[Lemma 3.2]{hhs}.

Summing the loads on a pattern \(A\) of size \(r\) gives
\begin{equation}\label{eq:edge}
21r\ \ge\sum_{x\in A}\ell_x
= W+(r-1)w_A+\sum_{B\ne A} (|A\cap B|-1)w_B .
\end{equation}
All summands on the right are nonnegative. Thus every pattern has size
at least four, each size-four pattern weighs less than 5, and each
size-five pattern weighs less than 9. Write \(\F\) for the four-element
patterns, \(t=|\F|\), and \(\G\) for the remaining patterns. Set
\[
 H=\sum_{B\in\G}w_B,\qquad E=\sum_B(|B|-4)w_B,\qquad
 S=\sum_{x\in V}(21-\ell_x)=294-4W-E .
\]
Then \(H\le E<18\). Since \(W-H>51\), we have \(11\le t\le14\).
Every vertex lies in a member of \(\F\), since
\(\ell_x\ge21-S=4W+E-273>3+E>H\).
Define the positive deficits \(\delta_A=5-w_A\) for \(A\in\F\) and
\[
 D=\sum_{A\in\F}\delta_A=5t-W+H .
\]
Equation~\eqref{eq:edge}, applied to any \(A\in\F\), gives
\begin{equation}\label{eq:repeats}
 \sum_{B\ne A} (|A\cap B|-1)w_B < 3\delta_A .
\end{equation}

The same estimates also explain what is needed to beat the
projective plane. With loads bounded by one and mass \(M>13/4\),
the support reduction gives at most fourteen patterns. Summing loads
on an individual pattern shows that its size is at least four and that
each four-set weighs at most \((4-M)/3<1/4\). The weight on larger
patterns is at most \(14-4M<1\), so the four-sets carry more than
\(9/4\) and number at least ten.

\begin{lemma}\label{lem:small}
\(D\ge4\).
\end{lemma}
\begin{proof}
Let \(d_x\) be the number of members of \(\F\) containing \(x\).
If \(D<4\), then \(d_x\le4\) for every \(x\), because five incident
members would weigh at least \(25-D>21\). For each \(A\in\F\),
\(\sum_{x\in A} d_x\ge t+3\) by counting the intersection with every
member, including \(A\) itself. Since every vertex lies in some \(A\)
and the other three degrees sum to at most 12, \(d_x\ge t-9\).
Also \(\sum_xd_x=4t\). Double counting now gives
\[
(t-9)(56-4t)\ \le\ \sum_x d_x(4-d_x)
=\sum_{A\in\F}\sum_{x\in A}(4-d_x)\ \le\ t(13-t).
\]
For \(t=11,12,13,14\) the left and right sides are respectively
\((24,22),(24,12),(16,0),(0,-14)\). Each contradicts the inequality.
\end{proof}

\begin{lemma}\label{lem:outside}
Any vertex outside all members of \(\G\) lies in at least five members
of \(\F\). If \(\G\ne\varnothing\) and \(O\) is a set of such vertices,
then \(4|O|\le 3D\).
\end{lemma}
\begin{proof}
Let \(x\) be such a vertex and \(k=d_x\ge1\). For each incident four-set
\(A\), summing the loads on \(A\setminus\{x\}\) gives
\(63\ge W-\ell_x+3w_A\), since every pattern missing \(x\) meets
\(A\) elsewhere. Summing over the \(k\) choices of \(A\) gives
\(63k\ge kW+(3-k)\ell_x\). For \(k\le3\) this implies \(W\le63\).
For \(k=4\) it implies \(4W\le252+21\). Thus \(k\ge5\) and the
deficit of the incident four-sets is \(5k-\ell_x\ge4\).
Each four-set meets a member of \(\G\), so it has at most three vertices
in \(O\). Summing the deficits on \(O\) proves the claim.
\end{proof}

\section{Reduction to two five-element patterns}
By Lemma~\ref{lem:small}, \(5t=W-H+D>69-18+4=55\).
Thus \(t\ge12\) and \(|\G|\le2\). If \(\G\) were empty then
\(D=5t-W<1\). If \(\G=\{Q\}\), then \(t\le13\), so
\(w_Q=D+W-5t>8\). Since \(E=(|Q|-4)w_Q<18\), \(|Q|\le6\).
For \(|Q|=5\), equation~\eqref{eq:edge} gives \(w_Q<9\); for
\(|Q|=6\), the same bound follows from \(E<18\). Consequently \(D<5\).
At least eight vertices lie outside \(Q\), contradicting
Lemma~\ref{lem:outside}.

We have \(t=12\) and two members of \(\G\), of total weight
\(H=W-60+D>13\). At least one, \(Q\), has size five, since otherwise
\(E\ge2H\). Suppose the other, \(R\), has size at least six and weight
\(g\). Then \(g=H-w_Q>4\). If \(|R|\ge7\), then \(E\ge H+2g>21\).
If \(|R|=6\), then \(E=H+g=2H-w_Q<18\), hence \(H<27/2\) and
\(D<9/2\). As \(Q\) meets \(R\), at least four vertices are outside
both, which would give \(16\le 3D<27/2\).
It remains to consider two five-element patterns \(Q_1,Q_2\).

\section{Two five-element patterns}
Let \(I=Q_1\cap Q_2\), \(j=|I|\ge1\), and \(O=V\setminus(Q_1\cup Q_2)\),
so \(|O|=4+j\). Now \(E=H\) and \(D=60-W+H<9\).
Lemma~\ref{lem:outside} gives \(4(4+j)\le3D<27\), whence \(j\le2\).
Write \(m\) for the number of four-sets meeting \(I\) and \(D_I\)
for their total deficit. A four-set missing \(I\) meets the two
five-sets in different vertices and has at most two vertices in \(O\).
A four-set meeting \(I\) has at most three. Consequently
\begin{equation}\label{eq:outside}
 5|O|\le 24+m,\qquad 4|O|\le2D+D_I .
\end{equation}

\emph{Two common vertices.}
Here \(|O|=6\), \(D\ge8\) and \(H=W-60+D>17\). By
\eqref{eq:outside}, \(m\ge6\). The four-sets meeting \(I\) weigh at
least \(30-D>21\). But the two vertices in \(I\) can carry at most
\(2(21-H)<8\) from these four-sets.

\emph{One common vertex.}
Write \(I=\{v\}\) and \(s_v=21-\ell_v\). Then \(0\le s_v\le S\) and
\(D_I=5m-21+H+s_v\). If \(m\ge3\), the load at \(v\) is at least
\(H+15-D=W-45>24\). If \(m\le1\), the second bound in
\eqref{eq:outside} gives
\[
20\le2D+D_I\le104-2W+3H+s_v
 <20-3S+s_v\le20,
\]
where \(H=294-4W-S<18-S\). Hence \(m=2\), and the ten four-sets
missing \(v\) have total deficit \(D-D_I=71-W-s_v<2\).

Any two of those ten four-sets \(A,B\) meet in exactly one vertex. A
second common vertex would give \(3\delta_A+\delta_B>5\) and
\(3\delta_B+\delta_A>5\) by \eqref{eq:repeats}, so their total deficit
would exceed \(5/2\). Furthermore \(20\le3D\) gives \(H>47/3\).
Each five-set therefore weighs more than \(47/3-9=20/3\). A four-set
missing \(v\) has \(\delta_A<2\), so \eqref{eq:repeats} forces it to
meet each five-set exactly once. Its remaining two vertices are in \(O\).

Group the ten four-sets by their unique vertex in \(Q_1\setminus\{v\}\).
Within a group their outside pairs must be disjoint, since the four-sets
already share the vertex indexing their group. Five outside vertices
accommodate at most two disjoint pairs. The four groups can therefore
hold at most eight four-sets. This contradiction establishes
Theorem~\ref{thm:main}.

\section*{Prior work and verification}
Mills determined the relevant covering ratios through thirteen blocks
\cite{mills,hhs}. Guiduli and Kir\'aly distinguish the one-vertex
quantity studied here from their two-vertex parameter \(M_2\) \cite{gk}.
F\"uredi's uniform fractional matching bound \cite{furedi} gives
\(\nu^*\le 3+1/4\) for an intersecting 4-uniform hypergraph.
The order-three projective plane attains that value on thirteen vertices.
Our lower bound allows patterns of different sizes, as does Barker's
construction. The La Jolla archive credits Barker with the
\((69,21,2)\) covering on 28 December 2008 \cite{barker}. Among the
716 pair-covering entries of size at most fourteen in version 1.2.1 of
that archive, \(21/69\) is the least \(k/v\). A finite table of
constructions alone cannot certify a lower bound for arbitrary weights.

We recovered the witness using floating-point linear programs to
evaluate local changes to supports initialized from the order-three
projective plane. With capacities four, two of three runs found the
support attaining mass \(92/7\). The mathematical proof uses the
displayed integer witness and exact inequalities. A separate Python
verifier expands the witness and checks every point pair.
Lean~4.24.0 checks the construction and the universal lower bound over
real weights. The accompanying sources expose the normalized result
as \texttt{L14.mills\_eq\_seven\_twentyThree}.

\paragraph{Contributions.} Troy Valiyaparampil posed the research
question, proposed testing \(4/13\), and requested formal verification.
GPT-6.2 Astra, working in OpenAI Codex, coordinated the research,
developed and checked the argument, and drafted the manuscript. GPT-6.2
Sol sessions contributed the search, proof arguments, source checks and
Lean modules.

\begin{thebibliography}{9}\small
\bibitem{mills} W.~H. Mills, Covering designs I. Coverings by a small
number of subsets, \emph{Ars Combinatoria} \textbf{8} (1979), 199--315.
\bibitem{furedi} Z. F\"uredi, Maximum degree and fractional matchings in
uniform hypergraphs, \emph{Combinatorica} \textbf{1} (1981), 155--162.
\href{https://doi.org/10.1007/BF02579271}{doi:10.1007/BF02579271}.
\bibitem{furedi90} Z. F\"uredi, Covering pairs by \(q^2+q+1\) sets,
\emph{Journal of Combinatorial Theory, Series A} \textbf{54} (1990),
248--271, Conjecture (6.3), p.~257.
\href{https://doi.org/10.1016/0097-3165(90)90034-T}{doi:10.1016/0097-3165(90)90034-T}.
\bibitem{furedi94} Z. F\"uredi, Intersecting designs from linear
programming and graphs of diameter two, \emph{Discrete Mathematics}
\textbf{127} (1994), 187--207, Problem~4.8 and Conjecture~4.9, p.~193.
\href{https://doi.org/10.1016/0012-365X(92)00478-A}{doi:10.1016/0012-365X(92)00478-A}.
\bibitem{gk} B. Guiduli and Z. Kir\'aly, On intersecting hypergraphs,
\emph{Discrete Mathematics} \textbf{182} (1998), 139--151.
\href{https://doi.org/10.1016/S0012-365X(97)00136-2}{doi:10.1016/S0012-365X(97)00136-2}.
\bibitem{hhs} J. Hall, D. Horsley and D.~R. Stinson, Bounds on data
limits for all-to-all comparison from combinatorial designs,
\emph{Designs, Codes and Cryptography} \textbf{92} (2024), 2863--2879.
\href{https://doi.org/10.1007/s10623-024-01418-6}{doi:10.1007/s10623-024-01418-6}.
\bibitem{barker} C. Barker, 14-block covering of 69 points by 21-subsets
(2008), in D.~M. Gordon, \emph{La Jolla Coverings Repository}, v1.2.1,
10 November 2026, entry \(C(69,21,2)\).
\href{https://doi.org/10.5281/zenodo.24119701}{doi:10.5281/zenodo.24119701}.
\end{thebibliography}
\end{document}
